%!TEX root = ../graduate-work.tex
\section{Постановка задачи}

Основным ограничением современных верификаторов является потребление
GPU-памяти: алгоритмы bound propagation (IBP, CROWN, $\alpha$-CROWN)
требуют хранения весовых матриц и матриц линейной релаксации
целиком на одном ускорителе, что делает невозможной верификацию
крупных моделей.

\textbf{Цель работы}~--- адаптировать методы параллелизма,
разработанные для обучения больших языковых моделей, к задаче
формальной верификации нейронных сетей, и экспериментально
исследовать компромиссы между экономией памяти, точностью
bounds и совместимостью с полной верификацией.

\textbf{Задачи работы:}
\begin{enumerate}
    \item Разобраться, насколько методы параллелизма из обучения LLM~---
    Tensor Parallelism, Pipeline Parallelism, FSDP~--- применимы к
    алгоритмам bound propagation и какие из них имеет смысл адаптировать.

    \item Реализовать \textbf{Tensor Parallelism} в форке
    \texttt{auto\_LiRPA}: шардирование весовых матриц и матриц
    линейной релаксации $A$ по нескольким GPU, интеграция
    с ONNX-моделями, проверка soundness, измерение пиковой памяти.

    \item Реализовать \textbf{FSDP} для bound propagation~---
    послойное AllGather/free весов без изменения логики вычислений~---
    и добиться побитовой идентичности bounds однопроцессорному режиму.

    \item Интегрировать \textbf{FSDP с полной верификацией}
    ($\beta$-CROWN + Branch-and-Bound); провести честное сравнение,
    устранив артефакт авто-расширения батча, чтобы оба режима
    обрабатывали идентичный workload.

    \item Распространить FSDP-шардирование на \textbf{свёрточные слои}
    (\texttt{BoundConv}) и проверить корректность на моделях ResNet
    из бенчмарка VNN-COMP.

    \item Провести серию экспериментов на бенчмарках VNN-COMP
    (MNIST-FC, ViT, CIFAR-100 ResNet) и выяснить, при каких условиях
    каждый из методов даёт реальную экономию памяти.
\end{enumerate}
